SAGE-CREH obtuvo el mayor cumplimiento estricto de la comparación principal:
23/24 (95.8\%), con 1765 tokens y 7.34 segundos medios. Frente a los cuatro
patrones originales redujo tokens un 47.2\% respecto al secuencial, 49.3\% al
paralelo, 61.1\% al jerárquico y 30.2\% al reflexivo. Son diferencias observadas;
los intervalos de diferencias de acierto estricto incluyen cero. La ventaja media
de tiempo frente al paralelo fue de solo 0.35 segundos, con intervalo que incluye
cero; tampoco se establece una ventaja de latencia frente al reflexivo.

La conclusión depende del contrato de salida. Con SQL normalizado, secuencial y
reflexivo lograron 24/24; SAGE y el agente único original, 23/24. Este último
consumió 1160 tokens y 4.62 segundos medios, menos que SAGE. En el frente de Pareto
de los seis diseños principales aparecen agente único y SAGE con la métrica
estricta; al usar la secundaria aparecen agente único y reflexivo. Son fronteras
de puntos estimados, no pruebas de dominancia poblacional. SAGE no es, por tanto,
la mejor elección bajo todas las definiciones de éxito.

El control SAGE sin delegación logró 21/24 con 1110 tokens: añadir coordinación
elevó el consumo un 59.1\% y produjo dos aciertos más en promedio de la muestra,
sin una diferencia de cumplimiento concluyente. Sus dos fallos SQL contienen
una clave extra \texttt{uncertainties} en el sobre final; esa forma no está cubierta
por la normalización secundaria predefinida, que solo envuelve cadenas SQL.
Por ello, tampoco esta columna separa todo posible error de formato del contenido.

La inspección de las propias trazas de SAGE identifica cinco mejoras frente a su
primer candidato: tres revisiones SQL de contrato y dos arbitrajes en dependencias
y selección. No hubo regresiones estrictas tras revisión. Se aceptaron trece
respuestas directas y cinco acuerdos paralelos sin integrador; hubo tres arbitrajes
y tres revisiones selectivas, para 38 llamadas en 24 tareas. Esto documenta
mecanismos concretos, sin convertir cada corrección de contrato en razonamiento
multiagente superior.

El único fallo final de SAGE ocurrió en T02, primera repetición: todas las sumas
monetarias coincidían, pero se devolvieron siete eventos cuando correspondían
ocho. El contrato era válido, no se declararon dudas y la vía directa lo aceptó.
Este caso muestra exactamente qué no detecta la compuerta de formato.

Seleccionar instrucciones redujo tokens un 40.5\% respecto a SAGE con diez skills
(1765 frente a 2966). No redujo el tiempo observado: 7.34 frente a 6.38 segundos,
con diferencia incierta. Las respuestas también cambiaron las rutas: 38 frente
a 34 llamadas. Esa ablación mide el efecto total de la carga bajo la misma
política, no un grafo ejecutado idéntico. En la pareja jerárquica, ambos hicieron
96 llamadas y cargar diez skills elevó tokens un 89.3\%, sin aumentar el total
de aciertos normalizados (22/24 en ambos).

Para este contrato estricto, SAGE es una alternativa adaptativa de consumo bajo
frente a los cuatro equipos fijos. Si el sistema admite normalización del
envoltorio SQL y prioriza el coste, los datos favorecen la referencia individual;
si prioriza el cumplimiento normalizado observado, favorecen al reflexivo frente
al secuencial por su menor consumo y tiempo medios. Esta decisión condicionada
es más fiel a la evidencia que proclamar una arquitectura universalmente ganadora.
